\documentclass[11pt]{article}
\usepackage[margin=0.78in]{geometry}
\usepackage{booktabs,amsmath}
\usepackage{lmodern}
\title{Staffing and prices at incumbent fast-food restaurants after New Jersey's 1992 minimum-wage increase}
\author{}
\date{}
\begin{document}
\maketitle
\begin{abstract}
On April 1, New Jersey raised its hourly minimum from \$4.25 to \$5.05. Eastern Pennsylvania retained the old floor. Across restaurants with paired staffing reports or confirmed permanent closure, staffing rose by 0.47 full-time equivalents per New Jersey location and fell by 2.28 in Pennsylvania. The difference of 2.75 favors New Jersey, although a single comparison state and one pre-policy wave leave the counterfactual uncertain. Starting wages and tax-inclusive meal prices also rose relatively in New Jersey. The surveys identify neither the effect on new openings nor the change in industry-wide employment.
\end{abstract}

\section{What the comparison measures}
The baseline survey interviewed managers at 331 New Jersey and 79 eastern Pennsylvania locations of Burger King, KFC, Roy Rogers and Wendy's during February 15--March 4. The follow-up covered November 5--December 31 of the same year. The unit of analysis is a baseline location, given equal weight, so an outlet with ten workers counts as much as one with thirty. The employment estimand is the short-run change per incumbent location attributable to the policy, including losses from permanent exit. The operational outcomes describe establishments reporting the relevant item in both waves.

For each outcome, form a within-restaurant change $D_i=Y_{i,2}-Y_{i,1}$ and estimate
\[
  \widehat\tau=\overline D_{\rm NJ}-\overline D_{\rm PA}.
\]
Differencing removes persistent store differences. The Pennsylvania change supplies the missing New Jersey trend under an unchanged floor. A causal interpretation requires comparable average untreated changes, no relevant anticipation by baseline or spillovers into the comparison market, and sufficiently representative observed changes. Both states were surveyed across the same seasons and in the same chains. Local demand, tourism, or chain-specific shocks could still differ. A single pre-policy observation cannot establish parallel trends, and the baseline questionnaire itself describes the scheduled increase.

The reported standard error for the simple contrast is
$\sqrt{s^2_{\rm NJ}/n_{\rm NJ}+s^2_{\rm PA}/n_{\rm PA}}$, using the within-state variance of restaurant changes. Intervals use 1.96 standard errors. They quantify restaurant-level variation under independence. Statewide shocks do not average out across restaurants. With one treated and one comparison state, these intervals cannot quantify policy-level uncertainty.

\newpage
\section{Where the employment difference comes from}
Managers reported full-time and part-time nonmanagement workers separately from managers and assistants. Define staffing as full-time workers plus managers and assistants plus half of part-time workers. This full-time-equivalent (FTE) proxy imposes a part-time weight because individual hours are unavailable. All three counts must be present. The primary calculation assigns zero follow-up staffing to confirmed permanent closures. It leaves temporary closures and refusals missing because employees might have been retained.

\begin{table}[h!]
\centering
\caption{Employment per baseline restaurant with usable data}
\begin{tabular}{lrrrr}
\toprule
 & Locations & Before & After & Change\\
\midrule
New Jersey &309&20.43&20.90&0.47\\
Eastern Pennsylvania &75&23.38&21.10&$-2.28$\\
\midrule
NJ minus PA &&$-2.95$&$-0.20$&2.75\\
\bottomrule
\end{tabular}
\end{table}

The primary difference is 2.75 FTEs (restaurant-level 95\% interval 0.12 to 5.38). Most of the contrast comes from Pennsylvania's decline. On these same records the estimated differences in full-time workers, managers, and part-time workers are respectively $+3.16$, $-0.07$ and $-0.68$ people. Counting each worker equally yields a relative headcount change of 2.41 (interval $-0.51$ to 5.34). Part-time weights of one-quarter and three-quarters yield 2.92 and 2.58 FTEs. The sign survives these choices, but the precision and the economic quantity measured depend on the weighting.

Of the original locations, 321 in New Jersey and 78 in Pennsylvania answered the follow-up. Five New Jersey establishments and one Pennsylvania establishment closed permanently (1.51\% and 1.27\% of the respective baseline samples). Four other New Jersey outlets were closed for renovation, highway construction or a mall fire, and one refused. Ten New Jersey and two Pennsylvania records lack complete baseline staffing. Twelve and two, respectively, lack usable follow-up staffing under the primary rule. Unrecorded changes could differ systematically from the changes observed.

Among responding stores with complete counts, the contrast is 2.72 (304 NJ, 74 PA). Assigning all four temporary closures zero follow-up staffing gives 2.51 (313 NJ, 75 PA; interval $-0.13$ to 5.15). Deleting any one primary-sample location gives estimates from 2.22 to 3.09. Regressing the primary staffing change on a New Jersey indicator, chain indicators and company ownership yields 2.78 (heteroskedasticity-robust HC1 standard error 1.34). This allows chains and ownership types to have different average changes. None of these checks rules out a common Pennsylvania downturn.

\section{Which stores faced the largest increase?}
Among outlets reporting a baseline starting wage, 286 of 314 in New Jersey paid below \$5.05. At follow-up, 283 of 318 reported exactly \$5.05 and one reported less. Table~\ref{tab:ops} uses paired responses to measure the relative starting-wage change. The \$0.50 relative gain is consistent with a binding floor. Many stores started above the old minimum and needed a raise of less than \$0.80 to reach the new one.

Within New Jersey, define the initial shortfall as $G_i=\max(0,5.05-w_{i,1})$ in dollars per hour. Regress each change on $G_i$, chain indicators and ownership. A ten-cent larger shortfall predicts a 10.5-cent larger starting-wage increase ($n=301$) and a 0.30-FTE larger staffing change ($n=293$; HC1 interval $-0.02$ to 0.63 per ten cents). The staffing gradient is imprecise. The common baseline wage enters both shortfall and wage change, so baseline reporting error can strengthen the wage gradient mechanically. Initially low-paying stores may also face different labor supply or demand trends. These gradients describe exposure within the treated state.

\newpage
\section{How restaurants adjusted}
\begin{table}[h!]
\centering
\caption{Paired changes in wages and operations per reporting outlet}
\label{tab:ops}
\small
\begin{tabular}{lrrrrr}
\toprule
Outcome (unit) & $n_{\rm NJ}/n_{\rm PA}$ & $\Delta$NJ & $\Delta$PA & Difference & SE\\
\midrule
Starting wage (\$/hour) &301/68&0.467&$-0.030$&0.497&0.051\\
Meal basket (\$) &287/68&0.067&$-0.034$&0.101&0.044\\
Meal basket ($100\,\Delta\log$ price) &287/68&2.04&$-0.88$&2.92&1.30\\
Daily hours open &320/78&0.001&0.109&$-0.108$&0.103\\
Total registers &306/76&$-0.042$&0.132&$-0.174$&0.108\\
Registers open at 11 a.m. &298/74&$-0.040$&$-0.216$&0.176&0.102\\
Any employee meal benefit (pp) &321/78&$-0.31$&0.00&$-0.31$&0.54\\
Free employee meal offered (pp) &321/78&4.36&1.28&3.08&4.68\\
Usual first raise (\$/hour) &241/55&$-0.010$&$-0.017$&0.008&0.019\\
\bottomrule
\end{tabular}
\end{table}

The basket adds the tax-inclusive prices of a medium soda, small fries and a regular hamburger (two pieces of chicken or equivalent at KFC). Its relative increase is about ten cents, or 2.92 log points on the log specification (interval 0.37 to 5.47). Controlling for chain and ownership yields 3.19 log points. Higher customer prices are consistent with passing on some labor costs. The data contain neither sales quantities nor total labor costs, and the price measure includes taxes, so a pass-through rate cannot be recovered. Operating hours barely moved in New Jersey. The register, meal-benefit and raise estimates give limited evidence of systematic offsets; item response varies, and a reported first raise describes only stores giving a usable amount.

\section{Interpretation and scope}
A binding floor can raise starting pay and reduce desired labor at fixed demand in a competitive market. With wage-setting power, higher pay could instead improve recruitment and support employment. Lower turnover, changed staffing mix and demand shifts could also affect measured counts. The positive relative FTE change and the price increase constrain the description of what happened, but the surveys cannot distinguish these mechanisms or measure productivity, profits, or hours actually worked.

New stores cannot enter a panel drawn before April. Estimating entry requires comparable opening records in both markets; measuring aggregate employment requires coverage of entrants, exits and other employers, for example a longitudinal establishment payroll census. Those analyses remain incomplete. The evidence concerns the sampled incumbent chains over the remainder of 1992.

\small\raggedright\paragraph{Reproduction.} Run \texttt{python3 analysis.py} with NumPy in the supplied directory to regenerate \texttt{results.json}. The script reads fixed-width positions from \texttt{input/codebook} and dotted entries as missing in \texttt{input/observations.dat}; the supplied blank questionnaires define the items. Each row already pairs both waves. Sheet number 407 appears for two different state--chain combinations, so both rows are retained. One second-wave date has year code 91; excluding that record gives 2.79 FTEs. Survey windows and statutory rates follow the supplied task.
\end{document}
